\section*{अध्याय \ref{ch:g12:internal-heat} --- पृथ्वी की आंतरिक ऊष्मा}

\begin{solution}{exo:g12:internal-heat:1}
(a) \omterm{def:g12:internal-heat:heat-flow}{ऊष्मा-प्रवाह}; (b) \omterm{def:g12:internal-heat:heat-production}{रेडियोजनित ऊष्मा-उत्पादन}; (c) \omterm{def:g12:internal-heat:secular-cooling}{दीर्घकालिक शीतलन};
(d) \omterm{def:g12:internal-heat:convection}{ऊष्मीय संवहन}; (e) \omterm{def:g12:internal-heat:geothermal}{भूतापीय ऊर्जा}।
\end{solution}

\begin{solution}{exo:g12:internal-heat:2}
(1) \qty{0.0916}{W/m^2}। (2) \qty{4.67e13}{W}। (3) \qty{0.029}{K/m}: एक
डिग्री सेल्सियस का अंतर एक केल्विन का अंतर है।
\end{solution}

\begin{solution}{exo:g12:internal-heat:3}
थोरियम-232, यूरेनियम-238, पोटैशियम-40 और यूरेनियम-235। कमलैंड केवल
थोरियम-232 और यूरेनियम-238 के भू-न्यूट्रिनो देखता है।
\end{solution}

\begin{solution}{exo:g12:internal-heat:4}
$490/\sqrt{1} = \qty{490}{mW/m^2}$; $490/2 = \qty{245}{mW/m^2}$;
$490/5 = \qty{98}{mW/m^2}$; $490/7 = \qty{70}{mW/m^2}$, लगभग महाद्वीपों का
माध्य।
\end{solution}

\begin{solution}{exo:g12:internal-heat:5}
सतह के निकट भूमि का तापमान दिनों और ऋतुओं का अनुसरण करता है;
केवल कुछ मीटर नीचे जाकर वह वर्ष के माध्य तापमान पर स्थिर रहता है, और उससे भी
गहरे जाकर ही वह नीचे से आती ऊष्मा दिखाता है।
खोदने से छिद्र की चट्टान और जल हिल जाते हैं और उनका तापमान बदल जाता है,
इसलिए तब तक प्रतीक्षा की जाती है जब तक छिद्र चट्टान के तापमान पर
न लौट आए।
\end{solution}

\begin{solution}{exo:g12:internal-heat:6}
(1) $(265 - 10)/9.1 \approx \qty{28.0}{\celsius}$ प्रति किलोमीटर, यानी
\qty{0.0280}{K/m}। (2) $q = 2.72 \times 0.0280 \approx \qty{0.0762}{W/m^2}$,
लगभग \qty{76}{mW/m^2}। (3) महाद्वीपों के माध्य,
\qty{70.9}{mW/m^2}, से थोड़ा ऊपर, लगभग \qty{7}{\percent}: एक सामान्य महाद्वीपीय मान।
\end{solution}

\begin{solution}{exo:g12:internal-heat:7}
महाद्वीप: $0.0709 \times \num{2.073e14} \approx \qty{1.47e13}{W} =
\qty{14.7}{TW}$। महासागर: $0.1054 \times \num{3.028e14} \approx
\qty{3.19e13}{W} = \qty{31.9}{TW}$। कुल लगभग \qty{46.6}{TW}, जिसमें से
$31.9/46.6 \approx \qty{68}{\percent}$, दो-तिहाई, महासागरीय
तली से निकलता है।
\end{solution}

\begin{solution}{exo:g12:internal-heat:8}
(1) $\num{1.88e-6} \times 8000 \approx \qty{0.0150}{W/m^2} = \qty{15.0}{mW/m^2}$।
(2) $15.0/70.9 \approx \qty{21}{\percent}$, \omterm{def:g12:internal-heat:heat-flow}{ऊष्मा-प्रवाह} का पाँचवाँ भाग।
\end{solution}

\begin{solution}{exo:g12:internal-heat:9}
$\sqrt{t} = 490/70.9 \approx 6.91$, इसलिए $t \approx 47.8$ मिलियन वर्ष: लगभग
48 मिलियन वर्ष से पुरानी महासागरीय तली औसत महाद्वीप से कम
ऊष्मा खोती है।
\end{solution}

\begin{solution}{exo:g12:internal-heat:10}
यूरेनियम-235: $2^{1.4/0.704} = 2^{1.99} \approx 4$ गुना अधिक। पोटैशियम-40:
$2^{2.5/1.248} = 2^{2.00} \approx 4$ गुना अधिक। दोनों दो अर्ध-आयु पीछे हैं।
\end{solution}

\begin{solution}{exo:g12:internal-heat:11}
(1) $10^2/10^{-6} = \qty{1e8}{s}$, लगभग 3 वर्ष। (2) $(10^5)^2/10^{-6} =
\qty{1e16}{s}$, लगभग 320 मिलियन वर्ष। (3) $(\num{2.891e6})^2/10^{-6}
\approx \qty{8.4e18}{s}$, लगभग 265 अरब वर्ष। समय दूरी के
वर्ग के अनुपात में बढ़ता है: दस गुना मोटा, सौ गुना अधिक समय।
\end{solution}

\begin{solution}{exo:g12:internal-heat:12}
हर एक को $2^{3/T}$ से गुणा कीजिए: थोरियम-232, $10.4 \times 2^{3/14.0} \approx
\qty{12.1}{TW}$; यूरेनियम-238, $9.6 \times 2^{3/4.468} \approx \qty{15.3}{TW}$;
पोटैशियम-40, $4.0 \times 2^{3/1.248} \approx \qty{21.2}{TW}$; यूरेनियम-235,
$0.42 \times 2^{3/0.704} \approx \qty{8.1}{TW}$। कुल लगभग \qty{56.7}{TW},
आज के \qty{24.4}{TW} के दुगुने से अधिक। दोनों अल्पायु समस्थानिक उसका
$(21.2 + 8.1)/56.7 \approx \qty{52}{\percent}$ देते हैं, जबकि
आज $4.42/24.4 \approx \qty{18}{\percent}$।
\end{solution}

\begin{solution}{exo:g12:internal-heat:13}
शक्ति प्रति सेकंड द्रव्यमान, विशिष्ट ऊष्मा और तापवृद्धि का
गुणनफल है: $\num{3e8} = 1640 \times 4180 \times \Delta T$, इसलिए $\Delta T =
\num{3e8}/(1640 \times 4180) \approx \qty{44}{K}$। औसत महाद्वीप प्रति वर्ग मीटर
\qty{0.0709}{W} खोता है, इसलिए \qty{300}{MW},
$\num{3e8}/0.0709 \approx \qty{4.2e9}{m^2}$, यानी लगभग \qty{4200}{km^2}, की प्राकृतिक ऊष्मा-हानि है:
बिजलीघर वह ऊष्मा एकत्र करता है जिसे पृथ्वी स्वाभाविक रूप से हज़ारों
वर्ग किलोमीटर पर छोड़ती है।
\end{solution}

\begin{solution}{exo:g12:internal-heat:14}
(1) 100 मिलियन वर्ष, 260 अरब वर्ष से $\num{2.6e11}/\num{1e8} = 2600$ गुना
छोटा है; 200 मिलियन वर्ष 1300 गुना छोटा।
(2) संवहन ऊष्मा को चलती चट्टान के साथ ले जाता है: प्रावार को पार करने का
समय उसकी मोटाई को चाल से भाग देने पर मिलता है, प्रति वर्ष कुछ सेंटीमीटर पर $\qty{2891}{km}$,
लगभग सौ मिलियन वर्ष। चालन के पास ऐसा कोई
वाहक नहीं: उसका समय दूरी के वर्ग के अनुपात में बढ़ता है, और हज़ारों
किलोमीटर पर वह विशाल हो जाता है।
\end{solution}

\begin{solution}{exo:g12:internal-heat:15}
(1) $10/22.3 \approx \qty{45}{\percent}$; शेष स्वयं प्रावार का शीतलन है,
जो अपनी संचित ऊष्मा खोता है। (2) सीमा त्रिज्या $6371 - 2891 =
\qty{3480}{km}$ पर है; उसका क्षेत्रफल $4\pi (\num{3.48e6})^2 \approx
\qty{1.52e14}{m^2}$ है, और $\num{1e13}/\num{1.52e14} \approx
\qty{0.066}{W/m^2}$, लगभग \qty{66}{mW/m^2}, सतह पर के माध्य का लगभग
\qty{72}{\percent}। (3) ऊष्मा सीमा को चालन द्वारा पार करके
प्रावार की सबसे निचली परत में जाती है; गरम होकर यह चट्टान अपने ऊपर के प्रावार से
हल्की हो जाती है और ऊपर उठ सकती है, ऊष्मा को ऊपर ले जाते हुए: माना जाता है कि
प्लूम गहराई में इसी तरह शुरू होते हैं।
\end{solution}

\begin{solution}{pb:g12:internal-heat:1}
\textbf{1.} $4\pi (\num{6.371e6})^2 \approx \qty{5.101e14}{m^2}$, और
$2.073 + 3.028 = 5.101$: हाँ।

\textbf{2.} $3.028/5.101 \approx \qty{59}{\percent}$, लगभग तीन-पाँचवाँ।

\textbf{3.} महासागरीय तली नई है: वह कटकों पर गरम बनती है और
दूर जाते हुए ठंडी होती है, और नया \omterm{def:g11:lithosphere-asthenosphere:lithosphere}{स्थलमंडल} बहुत ऊष्मा खोता है।

\textbf{4.} $490/\sqrt{4} = \qty{245}{mW/m^2}$।

\textbf{5.} $2.72 \times 0.026 \approx \qty{0.0707}{W/m^2} =
\qty{70.7}{mW/m^2}$: माध्य से ठीक नीचे, व्यावहारिक रूप से उसके बराबर।

\textbf{6.} $0.0709 \times \num{2.073e14} \approx \qty{14.7}{TW}$।

\textbf{7.} $0.1054 \times \num{3.028e14} \approx \qty{31.9}{TW}$।

\textbf{8.} $14.7 + 31.9 \approx \qty{46.6}{TW}$, प्रकाशित अनुमान का
\qty{47}{TW} (अध्याय का वैश्विक माध्य \qty{46.7}{TW} देता है, वह
मान जो आगे उपयोग किया गया है)।

\textbf{9.} $\num{4.66e13}/\num{5.101e14} \approx \qty{0.0914}{W/m^2}$,
लगभग \qty{91}{mW/m^2}।

\textbf{10.} $0.0243 \times \num{2.073e14} \approx \qty{5.0}{TW}$, महाद्वीपों से होकर
खोए गए \qty{14.7}{TW} का एक-तिहाई: \omterm{def:g12:continental-crust:continental-crust}{महाद्वीपीय भूपर्पटी} की
रेडियोधर्मिता उनके \omterm{def:g12:internal-heat:heat-flow}{ऊष्मा-प्रवाह} का एक बड़ा भाग है।

\textbf{11.} यूरेनियम के हर किलोग्राम के लिए पृथ्वी में \qty{3.9}{kg}
थोरियम है, जो $3.9 \times 26.3 \approx \qty{102.6}{\micro W}$ मुक्त करता है,
जबकि यूरेनियम-238 \qty{94.4}{\micro W}। इसलिए यूरेनियम-238
\qty{20.0}{TW} का $94.4/(94.4 + 102.6) \approx 0.479$ बनाता है: लगभग
\qty{9.6}{TW}। यूरेनियम-235: $9.59 \times 4.10/94.4 \approx \qty{0.42}{TW}$।

\textbf{12.} $20.0 + 0.42 + 4 \approx \qty{24.4}{TW}$।

\textbf{13.} ग्रेनाइट और \omterm{def:g12:continental-crust:continental-crust}{महाद्वीपीय भूपर्पटी} की दूसरी चट्टानें, प्रावार की अतिमैफ़िक
चट्टानों से लगभग बीस गुना अधिक समृद्ध।

\textbf{14.} अधिक: तब हर समस्थानिक अधिक प्रचुर था, हर अर्ध-आयु पीछे
दुगुना; 3 अरब वर्ष पहले कुल आज के दुगुने से
अधिक था।

\textbf{15.} $46.7 - 24.4 = \qty{22.3}{TW}$।

\textbf{16.} पृथ्वी के बनने के समय से संचित ऊष्मा से: प्रावार और क्रोड का
\omterm{def:g12:internal-heat:secular-cooling}{दीर्घकालिक शीतलन}।

\textbf{17.} $22.3/46.7 \approx \qty{48}{\percent}$।

\textbf{18.} $11/46.7 \approx \qty{24}{\percent}$ से $38/46.7 \approx
\qty{81}{\percent}$ तक: भू-न्यूट्रिनो इस परास को बहुत संकरा कर देते हैं।

\textbf{19.} \omterm{def:g12:internal-heat:convection}{प्रावार संवहन}, जो अधिकांश ऊष्मा ले जाता है; केवल अंतिम
दसियों किलोमीटर चालन द्वारा पार किए जाते हैं।

\textbf{20.} $24.4/46.7 \approx 0.52$: आज रेडियोधर्मिता पृथ्वी की ऊष्मा-हानि का
\textbf{लगभग आधा} देती है, और ग्रह का धीमा शीतलन
दूसरा आधा।
\end{solution}
